%%%%%%%%%%%%%%%%%%%%%%%%%%%%%%%%%%%%%%%%%%%%%%%%%%%%%%%%%%%%%%%%%%%%%%%%%%%%%%%%
%%%%%%%%%%%%%%%%%%   Vorlage für eine Abschlussarbeit   %%%%%%%%%%%%%%%%%%%%%%%%
%%%%%%%%%%%%%%%%%%%%%%%%%%%%%%%%%%%%%%%%%%%%%%%%%%%%%%%%%%%%%%%%%%%%%%%%%%%%%%%%

% Erstellt von Maximilian Nöthe, <maximilian.noethe@tu-dortmund.de>
% ausgelegt für lualatex und Biblatex mit biber

% Kompilieren mit
% lualatex dateiname.tex
% biber dateiname.bcf
% lualatex dateiname.tex
% lualatex dateiname.tex
% oder einfach mit:
% make

\documentclass[
  tucolor,
  BCOR=12mm,     % 12mm binding corrections, adjust to fit your binding
  parskip=half,  % new paragraphs start with half line vertical space
  open=any,      % chapters start on both odd and even pages
  cleardoublepage=plain,  % no header/footer on blank pages
]{tudothesis}


% Warning, if another latex run is needed
\usepackage[aux]{rerunfilecheck}

% just list chapters and sections in the toc, not subsections or smaller
\setcounter{tocdepth}{1}

%------------------------------------------------------------------------------
%------------------------------ Sprache und Schrift: --------------------------
%------------------------------------------------------------------------------
\usepackage{fontspec}
\defaultfontfeatures{Ligatures=TeX}  % -- becomes en-dash etc.

% german language
\usepackage{polyglossia}
\setdefaultlanguage{german}

% for english abstract and english titles in the toc
\setotherlanguages{english}

% intelligent quotation marks, language and nesting sensitive
\usepackage[autostyle]{csquotes}

% microtypographical features, makes the text look nicer on the small scale
\usepackage{microtype}

%------------------------------------------------------------------------------
%------------------------ Für die Matheumgebung--------------------------------
%------------------------------------------------------------------------------

\usepackage{amsmath}
\usepackage{amssymb}
\usepackage{mathtools}

% Enable Unicode-Math and follow the ISO-Standards for typesetting math
\usepackage[
  math-style=ISO,
  bold-style=ISO,
  sans-style=italic,
  nabla=upright,
  partial=upright,
]{unicode-math}
\setmathfont{Latin Modern Math}

% nice, small fracs for the text with \sfrac{}{}
\usepackage{xfrac}


%------------------------------------------------------------------------------
%---------------------------- Numbers and Units -------------------------------
%------------------------------------------------------------------------------

\usepackage[
  locale=DE,
  separate-uncertainty=true,
  per-mode=symbol-or-fraction,
]{siunitx}
\sisetup{math-micro=\text{µ},text-micro=µ}

%------------------------------------------------------------------------------
%-------------------------------- tables  -------------------------------------
%------------------------------------------------------------------------------

\usepackage{booktabs}       % stellt \toprule, \midrule, \bottomrule

%------------------------------------------------------------------------------
%-------------------------------- graphics -------------------------------------
%------------------------------------------------------------------------------

\usepackage{graphicx}
\usepackage{grffile}

% allow figures to be placed in the running text by default:
\usepackage{scrhack}
\usepackage{float}
\floatplacement{figure}{htbp}
\floatplacement{table}{htbp}

% keep figures and tables in the section
\usepackage[section, below]{placeins}


%------------------------------------------------------------------------------
%---------------------- customize list environments ---------------------------
%------------------------------------------------------------------------------

\usepackage{enumitem}

%------------------------------------------------------------------------------
%------------------------------ Bibliographie ---------------------------------
%------------------------------------------------------------------------------

\usepackage[
  backend=biber,   % use modern biber backend
  autolang=hyphen, % load hyphenation rules for if language of bibentry is not
                   % german, has to be loaded with \setotherlanguages
                   % in the references.bib use langid={en} for english sources
]{biblatex}
\addbibresource{references.bib}  % die Bibliographie einbinden
\DefineBibliographyStrings{german}{andothers = {{et\,al\adddot}}}

%------------------------------------------------------------------------------
%------------------------------ Sonstiges: ------------------------------------
%------------------------------------------------------------------------------

\usepackage[pdfusetitle,unicode,linkbordercolor=tugreen]{hyperref}
\usepackage{bookmark}
\usepackage[shortcuts]{extdash}

%------------------------------------------------------------------------------
%-------------------------    Angaben zur Arbeit   ----------------------------
%------------------------------------------------------------------------------

\author{Dag-Björn Hering}
\title{Floquet Theorie für den inversen Faraday Effekt}
\date{2017}
\birthplace{Castrop-Rauxel}
\chair{Lehrstuhl für Theoretische Physik I}
\division{Fakultät Physik}
\thesisclass{Bachelor of Science}
\submissiondate{31. September 2015}
\firstcorrector{Prof.~Dr.~Erstgutachter}
\secondcorrector{Prof.~Dr.~Zweitgutachter}

% tu logo on top of the titlepage
\titlehead{\includegraphics[height=1.5cm]{logos/tu-logo.pdf}}

\begin{document}
\frontmatter
% \input{content/hints.tex}
\maketitle

% Gutachterseite
\makecorrectorpage

% hier beginnt der Vorspann, nummeriert in römischen Zahlen
\thispagestyle{plain}

\section*{Kurzfassung}
Hier steht eine Kurzfassung der Arbeit in deutscher Sprache inklusive der Zusammenfassung der
Ergebnisse.
Zusammen mit der englischen Zusammenfassung muss sie auf diese Seite passen.

\section*{Abstract}
\begin{english}
The abstract is a short summary of the thesis in English, together with the German summary it has to fit on this page.
\end{english}

\tableofcontents

\mainmatter
% Hier beginnt der Inhalt mit Seite 1 in arabischen Ziffern
\chapter{Einleitung}
In dieser Arbeit wird versucht,
den inversen Faraday Effekt
in einem Bandisolator mit
Hilfe der Floquet Theorie
zu beschreiben.

Der Inverse Farady Effekt (IFE) stellt eine Möglichkeit bereit,
die Magnetisierung in einem Material
durch Bestrahlung zu beeinflussen,
somit bietet er eine interessante
Methode für die
Kontrolle von Magnetisierungsprozessen
 in Nanostrukturen.
Durch ihn werden beispielsweise, wie schon
in Experimenten gezeigt \cite{jackl},
Spin-Wellen in einem Material erzeugt.
Weiterhin bietet er interessante
Aspekte für die Informationstechnik \cite{hertel}).
Als Grundlage für die
theoretische Beschreibung des IFE dient hier
der Artikel "Theory of the inverse Faraday effect
in metals" \cite{hertel} von Riccardo
Hertel, der sich mit dem IFE in Metallen
beschäftigt.

Die Floquet Theorie liefert eine Methode
zur nährungsweisen Berschreibung
getriebener quantenmechanischer Systeme.
Ein Vorteil dieser Theorie ist
die Berücksichtigung der Periodizität
der Störung auf jedem Nährungsschritt.
Ebenfalls ist das Fehlen von
säkularen Termen, die linear oder nicht
periodisch in der Zeit sind, von Vorteil
im Floquet Formalismus.\cite{haenggi}
Sie wird zum Beispiel in dem Artikel \cite{mentink}
dazu genutzt, die Austauschwechselwirkung
in Mott Isolatoren zu untersuchen.
Das Kapitel
"Driven Quantum Systems"
von Peter Hänggi aus dem Buch \cite{haenggi}
dient als Basis für die hier verwendete Floquet Theorie.


In dem Kapitel \ref{sec:theo} werden die
hier benötigten theoretischen Grundlagen
sowohl für den IFE als auch für die Floquet Theorie aufgeführt.
Anschließend wird in Kapitel \ref{sec:model} ein
Modell eines Bandisolators aufgestellt,
in dem der IFE untersucht wird.
Die durch das Modell gewonnenen Ergebnisse
werden in dem Kapitel \ref{sec:ergebnisse} dargelegt.
Abschließend werden in Kapitel \ref{sec:zusamm}
noch einmal die Vorgehensweise
und die gewonnenen Erkenntisse der vorherigen
Kapitel \ref{sec:model} und \ref{sec:ergebnisse}
zusammmengefasst und ein Ausblick
auf noch offene Aspekte gegeben.

%Im Rahmen dieser Bachelorarbeit wird sich mit der Thematik des Inversen Faraday Effektes (IFE) beschäftigt.
%Der IFE bietet eine Möglichkeit zur Kontrolle von ulta schnellen magnetisierungs Prozessen in Nanostrukturen

%sowie die kontrolle von Spinwellen in Materialien   blablabla ist.
%Welcher mit Hilfe der Floquet-Theorie untersucht werden soll.





% hier sollten Verweise auf die Paper stehen wieso inverser Farady effekt interessant
%und  wieso das Floquet theorem von Nutzen ist.

%
%
% \section{Inverser Faraday Effekt}
% Der Inverse Faraday Effekt beschreibt die Enstehung eines Kreisstromes in Materie, welcher
% eine Magnetisierung in dem Material hervorruft, bei Bestrahung mit zirkularpolariesiertem Licht.
%
% Der Inverse Faraday Effekt beschreibt die Enstehung einer Magnetisierung in Materie, die mit zirkularpolariesiertem Licht bestrahlt wird.
% Dabei wird davon ausgegangen, dass die Wechselwirkung primär zwischen dem E-felder der Welle und den Elektronen statt findet.
% Es kann durch mehrere Überlegungen gezeigt werden, dass der Strom, der durch das zirkularpolarisierte Licht erzeugt wird, eine quadratische Abhängigkeit
% zu der E-feld Amplitude besitzt.

\chapter{Theoretische Grundlagen}
\label{sec:theo}
Dieses Kapitel beschäftigt sich mit den
theoretischen Grundlagen, die für diese Arbeit benötigt
werden. Zunächst werden
quantenmechanische Postulate genannt,
die hier Verwendung finden.
Darauffolgend wird eine Methode dargelegt,
den IFE in Metallen
zu beschreiben. Schließlich wird versucht,
die gewonnenen Erkenntnisse über den
IFE in Metallen auf Isolatoren zu übertragen.
Es wird eine Einführung in
die Floquet Theorie gegeben, mit der es möglich
ist, periodische quantenmechanische Systeme
zu beschreiben. Abschließend
wird eine numerische Methode vorgestellt,
um die in der Floquet Theorie beschriebenen
Größen zu berechnen.
In den folgenden Gleichungen werden
\begin{align}
   \symup{e}=\hbar=\text{m}_\symup{e}=\frac{1}{4\pi\epsilon_0}=1
\end{align}
gesetzt.




\section{Postulate der Quantenmechanik}
In der Quantenmechanik wird
ein Zustand eines Systems durch seine Wellenfunktion $\ket{\Psi}$
beschrieben. Die zeitliche Entwicklung solch
eines quantenmechanischen Systems ist durch die
zeitabhängige Schrödingergleichung
\begin{align}
\mathrm{i} \frac{\partial}{\partial t}\ket{{\Psi(t)}}=  H(t) \ket{\Psi(t)} \label{eqn:schrodinger}
\end{align}
mit dem Hamiltonoperator $H(t)$
gegeben.
Für ein zeitunabhängigen Hamiltonoperator wird durch
Separation der Variablen
\begin{align}
  \ket{\Psi(t)}=f(t)\ket{\Psi}
\end{align}
die stationäre Schrödingergleichung
\begin{align}
H \ket{\Psi}_n=E_n\ket{\Psi}_n \label{eqn:stationSG}
\end{align}
mit den Eigenenergien $E_n$ und dem
Eigenzuständen $\ket{\Psi}_n$
hergeleitet.
% Für den zeitabhängigen Teil
% \begin{align}
% \symup{i}\frac{\partial}{\partial t}f(t)=Ef(t)
% \end{align}
% er ... mach ich später
Die Wahrscheinlichkeit $P(t)$ für eine
Wellenfunktion $\ket{\Psi(t)}$
in einem
Zustand $\ket{\phi}$ ist
durch
\begin{align}
  P(t)=\lvert\braket{\phi\vert\Psi}\rvert^2
\end{align}
gegeben.
Observablen des Systems werden durch hermitesche
Operatoren dargestellt.
Der Erwartungswert $\braket{O}$ solch einer Observablen zu
dem Operator $O$
ist durch
\begin{align}
\braket{O}=\braket{\Psi\lvert O\rvert\Psi}
\end{align}
definiert.
\cite{schwabl}


\section{Inverser Faraday Effekt}
\label{sec:inverfaraday}
Der IFE beschreibt die Entstehung eines Kreisstromes
in Materie durch Bestrahlung mit zirkular
 polarisiertem Licht.
Die dabei entstehende Magnetisierung
in der Materie bleibt gegebenenfalls
auch nach Bestrahlung permanent erhalten.
Hier wird zunächst eine Herleitung des
IFE in Metallen
dargelegt und im Anschluss versucht,
die Ergebnisse auf einen Isolator
zu übertragen.
Es wird davon ausgegangen, dass
die Wechselwirkung primär zwischen dem
oszillierenden elektrischen Feld der Welle
\begin{align}
  E(x,t)=\hat{E}\exp{\left(\symup{i}kr-\symup{i}\omega t\right)}
\end{align}
und den Elektronen stattfindet.
Des Weiteren werden ein wechselwirkungsfreies Elektronenplasma angenommen und quantenmechanische Effekte vernachlässigt.
Für die Bestimmung der Stromdichte $j$,
die durch zirkular polarisiertes Licht mit der Frequenz $\omega$ und
der Amplitude $\hat{E}$ des elektrischen Feldes entsteht, dient die Kontinuitätsgleichung
\begin{align}
  \frac{\partial}{\partial t}n +\nabla(nv) \label{eqn:konti}
\end{align}
mit der Elektronendichte $n$ und der Geschwindigkeit $v$
als Ausgangspunkt.
Wobei die Geschwindigkeit $v$ über den Zusammenhang
\begin{align}
  j=env=\sigma E
\end{align}
mit der Stromdichte $j$, dem elektrischen Feld $E$ und
der Leitfähigkeit des bestrahlten Materials $\sigma$
verknüpft ist.
Eine Größe $a$ lässt sich unter Betrachtung von zwei unterschiedlichen Zeitskalen darstellen als
\begin{align}
  a=\braket{a}+\delta a
\end{align}
mit $\braket{a}$ dem zeitlichen Mittelwert der Größe und
$\delta a$ dem oszillierenden Anteil in Zeiteinheiten des elektrischen Feldes.
Durch Anwenden dieser Annahme auf die Größen der
Kontinuitätsgleichung \eqref{eqn:konti}
wird ein Ausdruck für die entstehende Magnetisierung $M$
\begin{align}
  \vec{M}=-\frac{\symup{i}}{4 \braket{n}\omega }\left(\sigma^*\hat{E}^* \times \sigma\hat{E}\right)\label{eqn:magnet}
\end{align}
hergeleitet.
Aus dem Zusammenhang
\begin{align}
\braket{j}=\nabla\times M \label{eqn:nabla}
\end{align}
ergibt sich die mittlere Stromdichte $\braket{j}$,
die für die Magnetisierung $M$ verantwortlich ist.
Für den Fall einer links bzw. rechts zirkular polarisierten Welle
mit der Ausbreitungsrichtung in $\vec{e}_z$ gilt
\begin{align}
  \hat E \times \hat{E}^*=\pm\symup{i}\lvert E\rvert^2\cdot \vec{e}_z.
\end{align}
Durch Einsetzen der Leitfähigkeit
für ein wechselwirkungsfreies Elektronenplasma
\begin{align}
\sigma=\frac{\symup{i}\braket{n}}{\omega}
\end{align}
in die Gleichung \eqref{eqn:magnet}
ergibt sich die Magnetisierung in einem Metall,
welches mit rechts bzw. links zirkular polarisiertem
Licht in $\vec{e}_z$-Richtung
bestrahlt wird, zu
\begin{align}
  M=\pm\frac{\symup{i}\braket{n}}{4\omega^3}\lvert\hat{E}\rvert^2 \vec{e}_z.
\end{align}
Es ist eine Antiproportionalität
zur dritten Potenz in $\omega$
und eine quadratische Abhängigkeit in $\hat{E}$
für die Magnetisierung und somit auch für den Strommittelwert
zu erkennen. \cite{hertel}

Um eine Voraussage zu den Abhängigkeiten, die
bei dem IFE in einem Bandisolator auftreten,
zu treffen, wird die Leitfähigkeit
eines Isolators
\begin{align}
  \sigma=\symup{i}\omega C
\end{align}
mit der materialabhängigen Konstante $C$
ebenfalls in die Gleichung \eqref{eqn:magnet}
eingesetzt.\cite{fließebach}
Für die Magnetisierung in einem Bandisolator, der mit
rechts bzw. links zirkular polarisiertem Licht bestrahlt wird,
ergibt sich
\begin{align}
  \vec{M}=\pm\frac{\omega C^2}{4\braket{n}} \lvert \hat{E}  \rvert^2 \vec{e}_z. \label{eqn:iso}
\end{align}
Die Gleichung \eqref{eqn:iso}
sagt somit für den IFE in Isolatoren
eine Magnetisierung
vorher, die
linear
in der Frequenz $\omega$ und
quadratisch in dem Betrag der Amplitude des elektrischen
Feldes $\lvert\hat{E}\vert$ ist.
% Dies gilt ebenfalls
% für den zeitlichen
% Strommittelwert, welcher sich aus
% der Gleichung \eqref{eqn:nabla} ergibt.



\newpage
\section{Floquet Theorie}
\label{sec:floquetheo}
In diesem Abschnitt werden die Grundlagen der
Floquet Theorie vorgestellt. Diese wird
dazu genutzt,
um ein in der Zeit periodisches quantenmechanisches System zu beschreiben.
Die zeitliche Entwicklung für so ein System
ist durch die zeitabhängige
Schrödingergleichung  \eqref{eqn:schrodinger}
gegeben. Wobei der Hamiltonian die Bedingungen
\begin{align}
  &H(t)=H(t+T),
\intertext{und}
  H(t)&=H_0+V(t)  &\text{mit}&   &V(t)=V(t+T)
\end{align}
erfüllen muss.
ER ist also periodisch
in der Zeit
und kann als Summe von einem
ungestörten Anteil $H_0$
und einem in der Zeit periodischen
Potential $V(t)$ geschrieben werden.
Das Floquet Theorem besagt, dass Zustände
 $\ket{\Psi_{\alpha}(t)}$ existieren,
welche Lösungen
von \eqref{eqn:schrodinger} sind, die
die Form
\begin{align}
\ket{\Psi(t)_\alpha}=\exp\left(-\mathrm{i}\epsilon_\alpha\right)\ket{\Phi_\alpha(t)}\label{eqn:psi_a}
\intertext{haben. Dabei ist $\ket{\Phi_\alpha(t)}$ eine sogenannte Floquet Mode und
unterliegt der periodischen Bedingung}
\ket{\Phi_\alpha(t)}=\ket{\Phi_\alpha(t+T)}.
\end{align}
Der Parameter $\epsilon_\alpha$ ist real und
wird auch charakteristischer Exponent oder
Floquet Exponent genannt.
Mit Hilfe von \eqref{eqn:psi_a} lässt sich
die Schrödingergleichung \eqref{eqn:schrodinger}
in die Form
\begin{align}
\mathcal{H}(t)\ket{\Phi_\alpha(t)}=\epsilon_\alpha \ket{\Phi_\alpha(t)} \label{eqn:floquetgl.}
\intertext{bringen, wobei der Operator }
  \mathcal{H}(t)=\left( H(t)-\mathrm{i}\frac{\partial}{\partial t} \right)
\end{align}
ist. Durch die Ähnlichkeit der Gleichung
\eqref{eqn:floquetgl.} zu der stationären
Schrödingergleichung
werden $\epsilon_\alpha$ und $\phi_\alpha$
auch als Quasieigenenergien und Quasieigenzustände,
im Folgendem abgekürzt mit QEE und QEZ,
bezeichnet.
Darüber hinaus ist
\begin{align}
  \ket{\Phi_{\alpha '}(t)}=\ket{\Phi_\alpha(t)}\exp(\mathrm{i}n\omega t) \equiv \ket{\Phi_{\alpha n}(t)}
\end{align}
ebenfalls eine Lösung von \eqref{eqn:floquetgl.},
wobei $n \in \mathbb{Z} / \{0 \} $
und $\omega$ die Frequenz des Potentials ist.
Dies hat eine Änderung der QEE in der Form
\begin{align}
    \epsilon_\alpha \rightarrow \epsilon_{\alpha '}=\epsilon_\alpha+n\omega\equiv\epsilon_{\alpha n} \label{eqn:epsilon_n}
\end{align}
zufolge.
Durch diese Feststellung ist es
möglich, alle QEE
für unterschiedliche $\alpha$
in der ersten Brillouin-Zone darzustellen,
da sich alle weiteren
QEE über \eqref{eqn:epsilon_n}
berechnen lassen.
Die QEZ unterliegen der Orthogonaliätsbedingung
\begin{align}
  \braket{\braket{\Phi_{\alpha' }(t)|\Phi_{\beta'}(t)}}\equiv \frac{1}{T} \int_0^T\mathrm{d}t
  \braket{\Phi_{\alpha' }(t)|\Phi_{\beta'}(t)}=\delta_{\alpha'\beta'}=\delta_{\alpha\beta}\delta_{nm}. \label{eqn:ortho}
\end{align}
Um die zeitliche Entwicklung
eines Startzustandes $\ket{\Psi(0)}$
mit der Floquet Theorie durchzuführen,
wird
zunächst der Startzustand $\ket{\Psi(0)}$ durch
die Superposition der QEZ
\begin{align}
  \ket{\Psi(0)}=\sum_\alpha c_a\ket{\Phi_\alpha(0)}& &\text{mit}&  &c_\alpha=\braket{\Phi_\alpha|\Psi(0)} \label{eqn:super}
\end{align}
ausgedrückt.
Der Zeitpropagator der Floquet Theorie
\begin{align}
  K(t;0)=\sum_\alpha \exp\left(-\mathrm{i}\epsilon_\alpha t \right)\ket{\Phi_\alpha(t)}\bra{\Phi_\alpha(0)} \label{eqn:Propagator}
\end{align}
ermöglicht,
den Startzustand \eqref{eqn:super}
in der Zeit zu propagieren.
Die Anwendung dieses Propagators auf den Startzustand
\eqref{eqn:super}
liefert
\begin{align}
  \ket{\Psi(t)}=K(t;0) \ket{\Psi(0)}=\sum_\alpha c_\alpha \exp\left(-\mathrm{i}\epsilon_\alpha t \right)\ket{\Phi_\alpha(t)}. \label{eqn:psi_t}
\end{align}
\cite{haenggi},\cite{dr}
% !!!!! in die Ergebnisse !!!!! Mit Hilfe dieses Propagators ist es unkomplizit in dem Floquet Formalismus, zeitlichgemittelte (read) Erwartungswerte von
% Operatoren $\hat O$ zu berechnen.
% Das zeitliche Mittel ist gegeben durch
% \begin{align}
%   \bar{\hat O}= \frac{1}{T}\int_0^T \braket{\Psi(t)|O|\Psi(t)}
% \intertext{Durch Einsetzen von \eqref{eqn:psi_t} und \eqref{eqn:fourier} ergibt sich nach Ausführung der Integration}
%  \bar{\hat O}= \sum_\alpha \lvert c_\alpha \rvert^2  \sum_{-\infty}^{\infty} c_\alpha^n(x)^\dag \hat{O} c_\alpha^n (x)^{\phantom{\dag}}.
% \end{align}
%
% Eine Herleitung dieser Formel ist im Anhang zu finden.??
\newpage
\section{Floquet Matrix Methode}
\label{sec:matrix}
Im Folgenden wird die numerische Methode
der Floquet Matrix,
welche die Berechnung der QEE
$\epsilon_{\alpha}$ und QEZ
$\ket{\Phi_\alpha}$  ermöglicht, vorgestellt.
Die QEZ $\Phi_\alpha$ lassen sich, da es
sich um in der Zeit periodische Zustände handelt,
als Fourierreihe
\begin{align}
  \ket{\Phi_{\alpha}(t)}&=\lim_{N\to\infty}\sum_{n=-N}^{N} \exp(\mathrm{i}n\omega t) \ket{c_{\alpha}^n} \label{eqn:fourier}
 \intertext{mit}
 \ket{c_\alpha^n}&=\sum_{k=1}^D c_{\alpha,k}^n \ket{\varphi_k} %phi beliebiges orthonomal set auf h_0
\end{align}
darstellen.
Hierbei ist $\ket{\varphi_k}$ ein beliebiges
orthonormales Set auf $H_0$ mit der Dimension $D$.
Durch Einsetzen von \eqref{eqn:fourier} in \eqref{eqn:floquetgl.}
ergibt sich die Definition der Matrixelemente
\begin{align}
  \braket{\bra{\varphi_j m}\mathcal{H}_\mathrm{F}\ket{\varphi_k n}}\equiv \braket{\varphi_j|H^{m-n}|\varphi_k} + n \omega \delta_{n,m}\delta_{j,k} \label{eqn:H_f}
\intertext{mit}
H^{m-n}=\frac{1}{T}\int_0^T \mathrm{d}t H(t) \exp\left(-\mathrm{i}(m-n)\omega t\right) \label{eqn:H_n_m}
\end{align}
mittels denen sich die Floquet Matrix $\mathcal{H}_\mathrm{F}$ aufstellen lässt.
Aus der Lösung der Eigenwertgleichung
\begin{align}
  \mathrm{det}|\mathcal{H}_\mathrm{F}-\epsilon\mathbb{1}|=0
\end{align}
folgen die QEE $\epsilon_{\alpha,n}$ und die Eigenvektoren $\ket{\epsilon_{\alpha,n}}$.
Diese QEE
genügen den periodischen Bedingungen aus der Gleichung \eqref{eqn:epsilon_n}.
Um aus den Eigenvektoren $\ket{\epsilon_{\alpha,n}}$
die QEZ  $\ket{\Phi_{\alpha}(t)}$
zu bestimmen, wird die Gleichung \ref{eqn:fourier}
verwendet.
Die vektorartigen Entwicklungskoeffizienten $\ket{c_\alpha^n}$ können den
Eigenvektoren der QEE, die sich in der 1. Brillouin-Zone
befinden, entnommen werden.
Dabei entsprechen die Komponenten
der Eigenvektoren, welche zu der
$n$-ten Fourier Mode gehören,
den Komponeten $c_{\alpha,k}^n$.
\cite{haenggi},\cite{dr}

\chapter{Ergebnisse}
\section{Problem}
Hamiltonian des Systems
\begin{align}
  H=J\sum_{i=1}^4 \left(c_{i+1}^\dag c_i + c_{i}^\dag c_{i+1}   +a_i c_i^\dag c_i +\epsilon_i c_i^\dag c_i\right)\\
  H=\underbrace{J\sum_{i=1}^4 \left(c_{i+1}^\dag c_i + c_{i}^\dag c_{i+1}c_i\right)   +\sum_{i=1}^4a_i c_i^\dag c_i}{H_0} +\underbrace{\sum_{i=1}^4\epsilon_i c_i^\dag}{V(t)}
\end{align}
Handel sich um einen Band-Isolator $=> a_1=a_3=-a$ und  $a_2=a_4=a$
\begin{align}
  H_0=\begin{pmatrix}
  -a& J& & J \\
  J& a &J & 0\\
  0& J& -a& J \\
  J& 0&  J & a
\end{pmatrix}
\end{align}


\section{Floquetmatrix}
Durch Anwendung von \eqref{eqn:H_n_m} ergibt sich:
\begin{align}
  H^{m-n}=H_0\delta_{m,n} -\frac{E}{2}\left(R_{-} \exp\left( \mathrm{i}\upphi\right)\delta_{n+1,m}   +  R_{+}\exp\left( -\mathrm{i}\upphi\right)\delta_{n-1,m}\right)\\
  R_{-}=
\begin{pmatrix}
  r_{1_x}-r_{1_y}& & & \\
  &r_{2_x}-r_{2_y} & & \\
  & & r_{3_x}-r_{3_y}& \\
  & & & r_{4_x}-r_{4_y}
\end{pmatrix}
\end{align}




\section{Abbildungen}

\chapter{Zusammenfassungund Ausblick}


\appendix
% Hier beginnt der Anhang, nummeriert in lateinischen Buchstaben
\chapter{Ein Anhangskapitel}

Hier könnte ein Anhang stehen, falls Sie z.B. Code, Konstruktionszeichnungen oder Ähnliches mit in die Arbeit bringen wollen. Im Normalfall stehen jedoch alle Ihre Resultate im Hauptteil der Bachelorarbeit und ein Anhang ist überflüssig.


\backmatter
\printbibliography

\cleardoublepage
\thispagestyle{empty}
\section*{Eidesstattliche Versicherung}
Ich versichere hiermit an Eides statt, dass ich die vorliegende Abschlussarbeit mit dem Titel \enquote{\thetitle} selbstständig und ohne unzulässige fremde Hilfe erbracht habe.
Ich habe keine anderen als die angegebenen Quellen und Hilfsmittel benutzt, sowie wörtliche und sinngemäße Zitate kenntlich gemacht. 
Die Arbeit hat in gleicher oder ähnlicher Form noch keiner Prüfungsbehörde vorgelegen.

\vspace*{1cm}\noindent
\begin{center}
  \begin{tabular}{@{}p{0.4\textwidth}@{\hspace{0.15\textwidth}}p{0.4\textwidth}@{}}
  \rule{\linewidth}{0.25pt}& \rule{\linewidth}{0.25pt}\\
  Ort, Datum & Unterschrift
  \end{tabular}
\end{center}

\subsection*{Belehrung}
Wer vorsätzlich gegen eine die Täuschung über Prüfungsleistungen betreffende Regelung einer Hochschulprüfungsordnung verstößt, handelt ordnungswidrig.
Die Ordnungswidrigkeit kann mit einer Geldbuße von bis zu \SI[round-mode=places, round-precision=2]{50000}{€} geahndet werden. 
Zuständige Verwaltungsbehörde für die Verfolgung und Ahndung von Ordnungswidrigkeiten ist der Kanzler/die Kanzlerin der Technischen Universität Dortmund. 
Im Falle eines mehrfachen oder sonstigen schwerwiegenden Täuschungsversuches kann der Prüfling zudem exmatrikuliert werden \mbox{(\S\,63 Abs. 5 Hochschulgesetz --HG--).}

Die Abgabe einer falschen Versicherung an Eides statt wird mit Freiheitsstrafe bis zu 3 Jahren oder mit Geldstrafe bestraft.

Die Technische Universität Dortmund wird ggf.\ elektronische Vergleichswerkzeuge (wie z.\,B.\ die Software \enquote{turnitin}) zur Überprüfung von Ordnungswidrigkeiten in Prüfungsverfahren nutzen. \\[\baselineskip]

\noindent Die oben stehende Belehrung habe ich zur Kenntnis genommen.\\[1cm]
\begin{center}
\begin{tabular}{@{}p{0.4\textwidth}@{\hspace{0.15\textwidth}}p{0.4\textwidth}@{}}
\rule{\linewidth}{0.25pt}& \rule{\linewidth}{0.25pt}\\
Ort, Datum & Unterschrift
\end{tabular}
\end{center}

\end{document}
